\subsection{Separable multisections and the Picard variety}
\label{subsec:positive-genus-separable-picard}

For a general hyperplane section, the classical comparison of Picard
identity components gives a finite unipotent restriction kernel in positive
characteristic; see the account in
\cite[Remark~5.8, pp.~39--41]{KleimanPicard}.
In the equality setting with the additional ruledness hypothesis,
the distinguished normalized cover identifies
the Picard variety of the surface and excludes algebraically trivial
kernel classes.

\begin{lemma}\label{lem:positive-selfintersection-separable}
Let $V$ be an integral projective surface over an algebraically
closed field $k$ of characteristic $p>0$, smooth on a neighbourhood
of a smooth integral Cartier curve $E$ with $E^2>0$.
Let $q:V\to C$ be a morphism to a smooth projective curve of genus
at least one. If $E$ dominates $C$ and $q$ is smooth at the generic
point of $E$, then $q|_E:E\to C$ is separable.
\end{lemma}

\begin{proof}
The conormal sequence on the smooth neighbourhood of $E$ is
short exact \cite[Tag~06AA]{Stacks}. Its three terms are locally
free, so dualizing gives
\[
 0\longrightarrow T_E\longrightarrow T_V|_E
 \longrightarrow\OO_E(E)\longrightarrow0.
\]
If $q|_E$ were inseparable, its differential would be zero
\cite[Tag~0C1C]{Stacks}. The tangent map
$T_V|_E\to(q|_E)^*T_C$ would therefore factor through
\[
 \OO_E(E)\longrightarrow(q|_E)^*T_C.
\]
This homomorphism is nonzero at the generic point, since $q$ is
smooth there. A nonzero map of line bundles on a smooth projective
curve has an effective zero divisor, so
\[
 E^2\leq\deg(q|_E)\bigl(2-2g(C)\bigr)\leq0,
\]
a contradiction. Smoothness of $V$ away from $E$ is not used.
\end{proof}

\begin{proposition}\label{prop:positive-genus-picard-variety}
Let $X$ be a smooth projective integral surface over an algebraically
closed field $k$ of characteristic $p>0$, let $Y\subset X$ be a smooth
integral ample divisor, and let $M\ne\OO_X$ have order $p$ with
$M|_Y\simeq\OO_Y$. Suppose
\[
 2g(Y)-2=pY^2.
\]
Let $\pi:N\to X$ be the distinguished normalized cover, and put
$D=\pi^*Y$. Assume that $N$ is birationally ruled, with a fixed
presentation $k(N)=k(C)(t)$ whose smooth projective ruling curve
$C$ has positive genus. Use the morphisms
\[
 h:N\longrightarrow C,\qquad f:X\longrightarrow B=C^{(p)}
\]
of Proposition~\ref{prop:positive-genus-picard-fibration}.
Then the following assertions hold.
\begin{enumerate}[label=\textup{(\roman*)}]
\item The finite morphisms $h_D:D\to C$ and $f_Y:Y\to B$ are
separable. Writing $d=Y^2$, $g=g(C)$, and
$b=\deg h_D=\deg f_Y$, one has
\[
 pd-2b(g-1)=\deg R_{D/C}=\deg R_{Y/B}\geq0,
\]
where $R$ denotes the ramification divisor defined by the different.
\item Every morphism from $X$ to an abelian variety factors uniquely
through $f$. After choosing compatible $k$-points, the induced
homomorphism
\[
 \operatorname{Alb}(X)\longrightarrow\operatorname{Alb}(B)
\]
is an isomorphism.
\item Pullback gives an isomorphism of abelian varieties
\[
 f^*:\Pic^0_{B/k}
 \xrightarrow{\ \sim\ }(\Pic^0_{X/k})_{\mathrm{red}}.
\]
\item The image of $M$ in $\operatorname{NS}(X)$ has order
exactly $p$. In particular, $M$ is not algebraically equivalent
to zero.
\item Put $F=k(B)$ and $\Gamma=X\times_B\Spec F$. The restriction
$M|_\Gamma$ is nontrivial and has order $p$.
\end{enumerate}
\end{proposition}

\begin{proof}
\emph{Separability.}
The conductor equality and
Proposition~\ref{prop:positive-genus-picard-fibration} give that
$D$ is smooth and ample, $N$ is smooth on a neighbourhood of $D$,
and $D^2=pd>0$. They also give a nonempty open $C_0\subset C$
with
\[
 N_{C_0}\simeq C_0\times\PP^1.
\]
Since $D$ is ample, its intersection with a general ruling fibre
is positive, so it dominates $C$. Its generic point lies in this
product open, where $h$ is smooth.
Lemma~\ref{lem:positive-selfintersection-separable} proves that
$h_D$ is separable.

The finite flat radicial map $\pi_D:D\to Y$ has degree $p$.
Both curves are smooth, so it identifies with relative Frobenius
through an isomorphism $Y\simeq D^{(p)}$
\cite[Tag~0CCY]{Stacks}. The relation
$f\pi=F_{C/k}h$ identifies $f_Y$ with the Frobenius twist of
$h_D$. Equivalently, the $p$-power isomorphism of fields carries
$k(D)/k(C)$ to $k(D)^p/k(C)^p$.
Hence $f_Y$ too is separable, and both maps have degree $b$.
Riemann--Hurwitz, including the wild different multiplicities,
gives
\[
 2g(D)-2=b\bigl(2g(C)-2\bigr)+\deg R_{D/C};
\]
see \cite[Tag~0C1B, Lemma~53.12.2 and its preceding proof]{Stacks}.
Since $g(D)=g(Y)$, this gives the displayed formula in (i);
the Frobenius-twisted formula gives the same ramification degree
for $Y\to B$.

\emph{Maps to abelian varieties.}
We recall the constancy argument for a morphism
$\PP^1\to Q$, where $Q$ is an abelian variety.
Translation trivializes $\Omega^1_{Q/k}$, and its invariant
differentials generate it everywhere. Their pullbacks to
$\PP^1$ vanish, since $H^0(\PP^1,\Omega^1_{\PP^1/k})=0$.
Thus the differential of the morphism is zero.
In characteristic $p$ one must still remove the possibility
of an inseparable nonconstant map.
The kernel of
$d:k(t)\to\Omega^1_{k(t)/k}$ is $k(t^p)$, as follows from
the basis $1,t,\ldots,t^{p-1}$ over $k(t^p)$.
Consequently the generic map factors through relative Frobenius
of $\PP^1$. The resulting rational map from its smooth
Frobenius twist to $Q$ extends to a morphism, because $Q$ is
proper \cite[Tag~0BXZ]{Stacks}.
Repeating this factorization, a nonconstant rational function
on the reduced image curve would pull back to an element of
$k(t^{p^n})$ for every $n$. Its nonzero order at a zero or pole
would be divisible by every $p^n$, which is impossible.
Therefore every morphism $\PP^1\to Q$ is constant.

Put $A=k(C)$, $K=k(X)$, and $L=k(N)$.
Proposition~\ref{prop:distinguished-ruled-cover} gives
\[
 A\cap K=A^p,\qquad AK=L,
\]
and a proper birational morphism $r:S\to N$, where $S$ is
obtained from $C\times\PP^1$ by finitely many point blowups.
Write $q=hr:S\to C$. After shrinking $C_0$, both
$S_{C_0}$ and $N_{C_0}$ identify with $C_0\times\PP^1$.

Let $\varphi:X\to Q$ be any morphism to an abelian variety, and set
$\Phi=\varphi\pi r:S\to Q$.
On every closed ruling fibre over $C_0$, the map $\Phi$ is
constant by the preceding argument.
Choose the constant section $[1:0]$ of the product ruling.
Restriction to that section defines a morphism
$\beta_0:C_0\to Q$.
The maps $\Phi$ and $\beta_0q$ agree at every closed point of
the reduced finite-type scheme $S_{C_0}$; their closed equalizer
is therefore all of $S_{C_0}$.
The map $\beta_0$ extends uniquely to $\beta:C\to Q$ by
normality of $C$ and properness of $Q$
\cite[Tag~0BXZ]{Stacks}.
Equality on the dense product open then gives
\[
 \varphi\pi r=\beta q\quad\hbox{on all of }S.
\]

If $\beta$ is constant, so is $\varphi$.
Otherwise let $W$ be its reduced image in $Q$.
It is an integral projective curve, also the reduced image of
$\varphi$; it need not be smooth or normal.
The two embeddings of $k(W)$ into $K$ and into $A$ agree in
$L=k(S)$. Their image therefore lies in
$A\cap K=A^p$.
The latter field is the image of $k(B)$ under relative Frobenius.
Thus the inclusion $k(W)\hookrightarrow k(B)$ gives a rational
map $B\dashrightarrow W$.
Since $B$ is a normal curve and $W$ is proper, it extends to a
morphism. Composing with $W\hookrightarrow Q$ gives
$\gamma:B\to Q$, with $\beta=\gamma F_{C/k}$.
The identities
$f\pi r=F_{C/k}q$ and $\varphi\pi r=\beta q$ now imply
\[
 \varphi=\gamma f.
\]
Indeed the equality holds on the generic point of $X$, since
$K\hookrightarrow L$ is injective, and then everywhere by
separatedness. The factorization is unique because $f$ is
surjective.

\emph{Albanese and Picard varieties.}
Choose $x_0\in X(k)$ and put $b_0=f(x_0)$.
Let $a_X:X\to\operatorname{Alb}(X)$ and
$a_B:B\to\operatorname{Alb}(B)$ be the pointed Albanese maps.
Their universal property and Picard duality are
\cite[Expos\'e~236, Corollary~3.2 and Theorem~3.3(iii),
pp.~236--238]{FGAPicard}.
Functoriality gives a homomorphism
\[
 u:\operatorname{Alb}(X)\to\operatorname{Alb}(B),
 \qquad ua_X=a_Bf.
\]
The factorization just proved writes $a_X=\gamma f$ for a
pointed morphism $\gamma:B\to\operatorname{Alb}(X)$.
Universality of $a_B$ supplies a homomorphism
$v:\operatorname{Alb}(B)\to\operatorname{Alb}(X)$ with
$va_B=\gamma$.
Then $uv\,a_Bf=a_Bf$, and surjectivity of $f$ gives
$uv\,a_B=a_B$.
Uniqueness in the Albanese universal property gives $uv=1$.
Likewise
$vu\,a_X=v\,a_Bf=\gamma f=a_X$ gives $vu=1$.
This proves (ii).

For completeness, the Picard schemes of $X$ and $B$ exist
because they are smooth projective geometrically integral
over $k$
\cite[Expos\'e~232, Theorem~3.1, pp.~148--149;
proof in Sections~4--5, pp.~150--155]{FGAPicard}.
Their identity components are proper
\cite[Expos\'e~236, Theorem~2.1(ii), pp.~231--232]{FGAPicard}.
Over the perfect field $k$, the canonical abelian subscheme
of the identity component is its reduction.
For $B$, the whole identity component is smooth:
$H^2(B,\OO_B)=0$, so
\cite[Expos\'e~236, Proposition~2.10(ii), pp.~235--236]{FGAPicard}
applies.
The functorial duality of the Albanese and Picard varieties
therefore identifies the dual of $u$ with
\[
 f^*:\Pic^0_{B/k}\longrightarrow(\Pic^0_{X/k})_{\mathrm{red}}.
\]
It is an isomorphism, proving (iii).
This argument makes no assumption that $\Pic^0_{X/k}$ itself
is reduced.

\emph{The torsion class.}
First, $f^*:\Pic(B)\to\Pic(X)$ is injective on line-bundle
classes. Indeed $f_*\OO_X=\OO_B$ by the proof of
Proposition~\ref{prop:positive-genus-picard-fibration}, so the
projection formula gives
$P\simeq f_*f^*P\simeq\OO_B$ whenever $f^*P$ is trivial.

We also need the following consequence of separability.
Let $j:V\to B$ be a finite separable morphism of smooth
projective integral curves, and let $P\ne\OO_B$ have order $p$.
Write $P=\OO_B(E)$ and choose $u\in k(B)^*$ with
$\operatorname{div}_B(u)=pE$.
The element $u$ is not a $p$th power in $k(B)$:
otherwise $u=v^p$ would give
$p(E-\operatorname{div}v)=0$, and the free abelian divisor
group would imply $P\simeq\OO_B$.
Thus $T^p-u$ defines the nontrivial purely inseparable
degree-$p$ generic extension of the associated $\mu_p$-torsor.

If $j^*P$ were trivial, a rational function $w$ on $V$ would
satisfy $\operatorname{div}_V(w)=j^*E$.
Then $j^*u/w^p$ would have divisor zero, and hence would lie
in $k^*$. Since $k$ is algebraically closed, this constant is
a $p$th power. It follows that $u$ becomes a $p$th power in
$k(V)$.
A finite separable extension cannot contain a nontrivial
purely inseparable extension of its base field, so this is
impossible. Therefore $j^*P$ is nontrivial.

Suppose now that $M$ is algebraically trivial.
A $k$-point of $\Pic^0_{X/k}$ is a $k$-point of its reduction.
Part (iii) gives $M=f^*P$ for some $P\in\Pic^0(B)(k)$.
Injectivity of $f^*$ makes $P$ nontrivial of order $p$.
But $f_Y:Y\to B$ is finite separable by (i), whereas
$f_Y^*P=M|_Y$ is trivial. The preceding argument is a
contradiction.
Thus $M$ is not algebraically trivial, and its nonzero image in
$\operatorname{NS}(X)=\Pic(X)/\Pic^0(X)(k)$ is killed by $p$.
Its order is therefore exactly $p$, proving (iv).

\emph{The generic fibre.}
The curve $\Gamma$ is proper, integral and regular: its local rings
are localizations of those of the smooth surface $X$.
Moreover $f_*\OO_X=\OO_B$ gives
$H^0(\Gamma,\OO_\Gamma)=F$. This equality can be checked by
localizing a finite affine cover and its compatibility conditions.
Geometric smoothness of $\Gamma$ is not required.

Suppose $M|_\Gamma$ were trivial. A nowhere-vanishing section
on $\Gamma$ is a rational section $e$ of $M$ on $X$.
Its divisor $E$ has no horizontal component, because the
horizontal prime divisors of $X$ are the closed points of
$\Gamma$, at which the section is a unit.
Thus $E$ is a vertical Cartier divisor and $M\simeq\OO_X(E)$;
see \cite[Tags~0BE6 and 0BE8]{Stacks} for rational-section
divisors and their detection on a normal scheme.
Under a fixed trivialization $M^p\simeq\OO_X$, the section
$e^p$ is a rational function $u$ with
\[
 \operatorname{div}_X(u)=pE.
\]
On the proper regular curve $\Gamma$ this function has no
zero or pole, so it is an invertible regular function and
$u\in F^*$. The detection of units at codimension-one points
is \cite[Tag~031T, Lemma~10.157.6(2)]{Stacks}.

Since $Y$ dominates $B$, it is not a component of $E$.
Choose a trivialization of $M|_Y$ and write the restricted
rational section as $w\in k(Y)^*$.
The two trivializations of $M^p|_Y$ differ by an invertible
global function on $Y$, hence by a constant $c\in k^*$.
Therefore $u|_Y=cw^p$. The constant is a $p$th power, and
the finite extension $k(Y)/F$ is separable by (i).
It cannot contain a nontrivial purely inseparable extension
of $F$, so $u=v^p$ for some $v\in F$.
It follows that
\[
 p\bigl(E-\operatorname{div}_X(v)\bigr)=0.
\]
The Weil divisor group is free abelian on the prime divisors.
Hence $E$ is principal and $M$ is trivial, a contradiction.
This proves (v). The argument allows arbitrary vertical
component coefficients, including those of multiple fibres;
it never assumes that $E$ is pulled back from $B$.
\end{proof}

The proposition leaves a specific torsion obstruction in the
positive-genus branch: a kernel class must have nonzero order-$p$
image in the N\'eron--Severi group and a nontrivial order-$p$
restriction to the regular generic fibre. Its geometric
normalization being rational does not make the Picard group
of that regular curve over $k(B)$ torsion-free.
It does not assert that every torsion line bundle descends to the
ruling base. When $p\geq5$, the ample equality formula
\eqref{eq:normalized-smallest-equality-euler}, together with
$\chi(\OO_N)=1-g(C)$, gives $g(C)=1+pn$ for an integer $n\geq0$.
Part~(i) then gives $d\geq2bn$ without requiring very ampleness.
If $Y$ is very ample, the additional bound $b\geq p+2$ follows from
Proposition~\ref{prop:positive-genus-picard-fibration}.
